\documentclass[11pt]{article}
\usepackage{amsmath, amssymb, amsthm}
\usepackage[margin=1in]{geometry}

\newtheorem{theorem}{Theorem}
\newtheorem{corollary}{Corollary}
\newtheorem{proposition}{Proposition}
\newtheorem{remark}{Remark}

\title{Norm Scaling in Arc-Cosine Kernel:\\Why Unconstrained Gradient Ascent on DA Utility Diverges}
\author{Investigation notes --- \texttt{validate\_utility/}}
\date{February 2026}

\begin{document}
\maketitle

\section{Setup}

We consider a sparse variational GP with a Poisson observation model:
\begin{align}
    r(x) &\sim \text{Poisson}(f(x)), \quad f(x) = \exp(A \lambda(x) + \lambda_0) \\
    \lambda &\sim \mathcal{GP}(0, K_\theta)
\end{align}
with variational approximation $q(\tilde{\lambda}) = \mathcal{N}(m, V)$ at inducing points $\{\tilde{z}_j\}_{j=1}^M$.

The approximate posterior at any point $x$ is:
\begin{align}
    \mu(x) &= k(x)^\top \tilde{K}^{-1} m \label{eq:post_mean}\\
    \sigma^2(x) &= K(x,x) + k(x)^\top \tilde{K}^{-1}(V - \tilde{K})\tilde{K}^{-1} k(x) \label{eq:post_var}
\end{align}
where $k(x) = [K(x, \tilde{z}_1), \ldots, K(x, \tilde{z}_M)]^\top$ and $\tilde{K} = K(\tilde{Z}, \tilde{Z})$.

The posterior covariance between two points $x, x'$ is:
\begin{equation}\label{eq:post_cov}
    \Sigma_q(x, x') = K(x, x') + k(x)^\top \tilde{K}^{-1}(V - \tilde{K})\tilde{K}^{-1} k(x')
\end{equation}

The distribution-aware (DA) utility of a candidate $x^*$, conditioned on a single target $x_t$, is:
\begin{equation}\label{eq:da_utility}
    U_{\text{DA}}(x^*) = H_{\text{marg}}(x^*) - H_{\text{cond}}(x^* \mid \lambda(x_t))
\end{equation}
where $H_{\text{marg}}$ is the marginal entropy of $r(x^*)$ under the GP posterior,
and $H_{\text{cond}}$ is the entropy after Gaussian conditioning on the observed $\lambda(x_t)$.

The conditioning formulas (standard Gaussian conditioning) are:
\begin{align}
    \mu_{\text{cond}}(x^*) &= \mu(x^*) + \frac{\Sigma_q(x^*, x_t)}{\Sigma_q(x_t, x_t)} \bigl(\lambda_t - \mu(x_t)\bigr) \label{eq:cond_mean}\\
    \sigma^2_{\text{cond}}(x^*) &= \sigma^2(x^*) - \frac{\Sigma_q(x^*, x_t)^2}{\Sigma_q(x_t, x_t)} \label{eq:cond_var}
\end{align}

The kernel is the first-order arc-cosine kernel:
\begin{equation}\label{eq:kernel}
    K(x, x') = \frac{1}{\pi} \sqrt{v_x \cdot v_{x'}} \cdot J(\theta)
\end{equation}
where
\begin{align}
    v_x &= x^\top C x + \sigma_0^2, \label{eq:vx}\\
    \cos\theta &= \frac{x^\top C x' + \sigma_0^2}{\sqrt{v_x \cdot v_{x'}}}, \label{eq:costheta}\\
    J(\theta) &= \sin\theta + (\pi - \theta)\cos\theta, \label{eq:J}
\end{align}
with $C$ the structured covariance matrix encoding the receptive field,
and $\sigma_0^2$ a bias variance parameter.

%----------------------------------------------------------------------
\section{Homogeneity of the Kernel ($\sigma_0^2 = 0$)}
%----------------------------------------------------------------------

\begin{theorem}[Cross-kernel proportionality]\label{thm:proportionality}
Let $\sigma_0^2 = 0$. Let $x^* = c \, x_t$ for scalar $c > 0$. Then for any point $z$:
\begin{equation}
    K(x^*, z) = c \cdot K(x_t, z)
\end{equation}
and consequently $k(x^*) = c \cdot k(x_t)$.
\end{theorem}

\begin{proof}
With $\sigma_0^2 = 0$, the quadratic forms become $v_{x^*} = c^2 (x_t^\top C x_t) = c^2 q$
where $q \triangleq x_t^\top C x_t$, and $v_{x_t} = q$.

For any point $z$ with $v_z = z^\top C z$:
\[
    \cos\theta(x^*, z) = \frac{(cx_t)^\top C z}{\sqrt{c^2 q \cdot v_z}}
    = \frac{c \, x_t^\top C z}{c \sqrt{q \cdot v_z}}
    = \frac{x_t^\top C z}{\sqrt{q \cdot v_z}}
    = \cos\theta(x_t, z).
\]
Since $\theta(x^*, z) = \theta(x_t, z)$, we have $J(\theta(x^*, z)) = J(\theta(x_t, z))$. Then:
\[
    K(x^*, z) = \frac{1}{\pi}\sqrt{c^2 q \cdot v_z} \cdot J(\theta(x_t, z))
    = c \cdot \frac{1}{\pi}\sqrt{q \cdot v_z} \cdot J(\theta(x_t, z))
    = c \cdot K(x_t, z).
\]
Applying this to each inducing point $z = \tilde{z}_j$ gives $k(x^*) = c \cdot k(x_t)$.
\end{proof}

\begin{corollary}[Perfect posterior correlation]\label{cor:rho}
Under the conditions of Theorem~\ref{thm:proportionality}, the posterior correlation
between $x^*$ and $x_t$ is exactly~1:
\begin{equation}
    \rho^2 \triangleq \frac{\Sigma_q(x^*, x_t)^2}{\Sigma_q(x^*, x^*) \cdot \Sigma_q(x_t, x_t)} = 1.
\end{equation}
\end{corollary}

\begin{proof}
Define $W \triangleq \tilde{K}^{-1}(V - \tilde{K})\tilde{K}^{-1}$.
From Theorem~\ref{thm:proportionality}, $k(x^*) = c \, k(x_t)$, and
$K(x^*, x_t) = c \, K(x_t, x_t) = c \, v_t$ (using $\theta(x^*, x_t) = 0$
when $x^* \parallel x_t$, so $J(0) = \pi$ and $K(x_t, x_t) = v_t$).

Substituting into~\eqref{eq:post_cov}:
\begin{align}
    \Sigma_q(x^*, x_t) &= c \, v_t + c \, k(x_t)^\top W \, k(x_t) = c \cdot \Sigma_q(x_t, x_t), \label{eq:cov_scale}\\
    \Sigma_q(x^*, x^*) &= c^2 v_t + c^2 k(x_t)^\top W \, k(x_t) = c^2 \cdot \Sigma_q(x_t, x_t). \label{eq:var_scale}
\end{align}
Therefore:
\[
    \rho^2 = \frac{[c \cdot \Sigma_q(x_t, x_t)]^2}{c^2 \cdot \Sigma_q(x_t, x_t) \cdot \Sigma_q(x_t, x_t)} = 1.
\]
\end{proof}

\begin{corollary}[Scaling of posterior moments]\label{cor:moments}
Under the conditions of Theorem~\ref{thm:proportionality}:
\begin{align}
    \mu(c \, x_t) &= c \cdot \mu(x_t), \label{eq:mu_scale}\\
    \sigma^2(c \, x_t) &= c^2 \cdot \sigma^2(x_t). \label{eq:var_scale2}
\end{align}
\end{corollary}

\begin{proof}
Direct substitution of $k(x^*) = c \, k(x_t)$ into~\eqref{eq:post_mean} and~\eqref{eq:post_var}.
\end{proof}

\begin{remark}[Effect of $\sigma_0^2 > 0$]\label{rem:sigma0}
When $\sigma_0^2 > 0$, the above results hold approximately.
The exact expressions are:
\begin{align}
    v_{cx_t} &= c^2 q + \sigma_0^2 \neq c^2 (q + \sigma_0^2) = c^2 v_t, \\
    \cos\theta(cx_t, z) &= \frac{c \, x_t^\top C z + \sigma_0^2}{\sqrt{(c^2 q + \sigma_0^2) v_z}}
    \neq \cos\theta(x_t, z) \quad \text{in general.}
\end{align}

The fractional error $|K(cx, z) - c \, K(x, z)| / |c \, K(x,z)|$ converges to a
\emph{constant} $O(\sigma_0^2/v_t)$ as $c \to \infty$, not to zero.
This is because $\sigma_0^2$ shifts the angle $\theta(cx_t, z)$ relative to
$\theta(x_t, z)$ by an amount that depends on $\sigma_0^2/v_t$, independent of $c$.

In our trained model, $v_t = q + \sigma_0^2 \approx 669$ with
$\sigma_0^2 \approx 0.98$, giving $\sigma_0^2 / v_t \approx 0.0015$.
Numerically verified (\texttt{verify\_conclusions.py}, test C7):
the mean fractional error in $K(cx_t, z_j)$ converges to
$\approx 0.6\%$ for large $c$ (slightly larger than $\sigma_0^2/v_t$
because both magnitude and angle corrections contribute).
The posterior correlation converges to $\rho^2 \approx 0.987$
(test C3), giving a residual variance ratio of $\approx 1.3\%$.
Conditioning remains very effective but not perfect.
\end{remark}

%----------------------------------------------------------------------
\section{Consequences for Conditioning}
%----------------------------------------------------------------------

From Corollary~\ref{cor:rho}, when $\rho^2 = 1$, the conditional variance~\eqref{eq:cond_var} becomes:
\begin{equation}
    \sigma^2_{\text{cond}}(x^*) = \sigma^2(x^*)(1 - \rho^2) = 0.
\end{equation}

The conditional mean~\eqref{eq:cond_mean} with $\lambda_t = \mu(x_t)$
(deterministic, i.e.\ \texttt{sample\_lambda=False}) and
using~\eqref{eq:cov_scale} becomes:
\begin{equation}
    \mu_{\text{cond}}(x^*) = c \, \mu(x_t) + c \cdot (\mu(x_t) - \mu(x_t)) = c \, \mu(x_t) = \mu(x^*).
\end{equation}

So in the exact parallel case ($\theta = 0$, $\sigma_0^2 = 0$):
\begin{itemize}
    \item The conditional mean equals the marginal mean: $\mu_{\text{cond}} = \mu$.
    \item The conditional variance is zero: $\sigma^2_{\text{cond}} = 0$.
    \item The conditional entropy reduces to the entropy of a single Poisson:
    \begin{equation}
        H_{\text{cond}} = H_{\text{Poisson}}\!\bigl(\exp(\mu_g)\bigr),
        \quad \mu_g \triangleq A \, \mu(x^*) + \lambda_0.
    \end{equation}
\end{itemize}

The DA utility~\eqref{eq:da_utility} becomes:
\begin{equation}\label{eq:utility_decomp}
    U_{\text{DA}}(c \, x_t) = H_{\text{marg}}(\mu_g, \sigma^2_g) - H_{\text{Poisson}}(\exp(\mu_g))
\end{equation}
where $\mu_g = A c \, \mu(x_t) + \lambda_0$ and $\sigma^2_g = A^2 c^2 \sigma^2(x_t)$.

%----------------------------------------------------------------------
\section{Why Utility Grows with $c$}
%----------------------------------------------------------------------

\begin{proposition}[Gaussian observation model: logarithmic growth]\label{prop:gaussian}
If the observation model were Gaussian, i.e.\
$r(x) \sim \mathcal{N}(\lambda(x), \tau^2)$ for fixed noise variance $\tau^2$,
then the DA utility grows as $\log(c)$ for $\rho = 1$ and $\sigma^2 = c^2 \sigma^2(x_t)$.
\end{proposition}

\begin{proof}
For Gaussian observations with GP prior variance $\sigma^2$ and noise $\tau^2$:
\begin{align}
    H_{\text{marg}} &= \tfrac{1}{2}\log\bigl(2\pi e (\sigma^2 + \tau^2)\bigr), \\
    H_{\text{cond}} &= \tfrac{1}{2}\log\bigl(2\pi e (\sigma^2_{\text{cond}} + \tau^2)\bigr)
    = \tfrac{1}{2}\log(2\pi e \, \tau^2) \quad [\text{since } \sigma^2_{\text{cond}} = 0].
\end{align}
Therefore:
\[
    U = \tfrac{1}{2}\log\!\bigl(1 + \sigma^2/\tau^2\bigr).
\]
Since $\sigma^2 = c^2 \sigma^2(x_t)$ (Corollary~\ref{cor:moments}), for large $c$:
\[
    U \approx \tfrac{1}{2}\log(c^2 \sigma^2(x_t)/\tau^2) = \log(c) + \text{const}.
\]
Growth is logarithmic and unbounded, but slow.
\end{proof}

\begin{remark}[What the Gaussian case shows]\label{rem:gaussian_growth}
The Gaussian case already exhibits unbounded growth --- norm scaling is not
specific to the Poisson model. However, because growth is only $\log(c)$,
the gradient $dU/dc \sim 1/c$ vanishes for large $c$, making gradient ascent
increasingly inefficient. A gradient-based optimizer would slow down but
never fully stop. The distinction with the Poisson case is the \emph{rate}
of growth, not its presence.
\end{remark}

For the Poisson-GP model, the marginal entropy involves:
\begin{equation}
    H_{\text{marg}}(\mu_g, \sigma^2_g) = -\sum_r p(r) \log p(r), \quad
    p(r) = \int \frac{e^{gr - e^g}}{r!} \, \mathcal{N}(g \mid \mu_g, \sigma^2_g) \, dg.
\end{equation}
This integral is computed via the Laplace approximation (see \texttt{\_diff\_laplace\_log\_probs}
in \texttt{utils.py}).

We do not have a closed-form expression for $H_{\text{marg}}(\mu_g, \sigma^2_g)$
in the Poisson-GP case. The following is a heuristic argument supported by
numerical evidence.

\paragraph{Heuristic argument.}
For large $\sigma^2_g$, the distribution $p(r)$ is a Poisson mixture
over log-normally distributed rates: $f = e^g$ with $g \sim \mathcal{N}(\mu_g, \sigma^2_g)$.

We claim $H_{\text{marg}} \geq H_{\text{Poisson}}(\mathbb{E}[f])$.
This is plausible because Poisson mixtures are overdispersed
($\text{Var}[R] > \mathbb{E}[R]$), hence ``more spread out'' than a single
Poisson at the same mean. However, a rigorous proof of this bound is not
provided here.  (Note: the Poisson does \emph{not} maximize entropy on
$\mathbb{N}_0$ under a mean constraint alone --- the geometric distribution
does. The Poisson maximizes entropy only under the joint constraint
$\text{Var} = \text{mean}$, so the max-entropy argument does not directly
apply.)  The bound is consistent with the numerical evidence
($\alpha \approx 1.9$ from \texttt{verify\_conclusions.py}).

With $\mathbb{E}[f] = e^{\mu_g + \sigma^2_g/2}$ and the large-$\lambda$
Poisson entropy $H_{\text{Poisson}}(\lambda) \approx \frac{1}{2}\log(2\pi e \lambda)$:
\[
    H_{\text{marg}} \geq \tfrac{1}{2}\log(2\pi e) + \tfrac{1}{2}(\mu_g + \sigma^2_g/2).
\]
The dominant term for large $c$ is $\frac{1}{4}\sigma^2_g = \frac{1}{4}A^2 c^2 \sigma^2(x_t)$.

Meanwhile, $H_{\text{cond}} = H_{\text{Poisson}}(\exp(\mu_g))$ depends only on $\mu_g$,
which grows as $c$ (not $c^2$). Therefore:
\[
    U_{\text{DA}} = H_{\text{marg}} - H_{\text{cond}} \gtrsim \tfrac{1}{4}\sigma^2_g \propto c^2,
\]
which is faster than the $\log(c)$ growth in the Gaussian case.

\textbf{Numerical calibration} (\texttt{verify\_conclusions.py}, test C4):
Fitting $U_{\text{DA}} \sim c^\alpha$ over the Laplace-valid range
($c \in \{2, 5, 10\}$, all with $\sum p(r) > 0.95$) gives
$\alpha \approx 1.9$, consistent with the predicted $c^2$ scaling.

\textbf{Caveat:} The Laplace approximation truncates at $r_{\max} = 100$
and breaks when $\sigma^2_g$ becomes large enough that the firing rate
distribution has significant mass at $r > r_{\max}$.
A useful pre-check is the ``Laplace safety score''
$z_{\text{safe}} = (\log r_{\max} - \mu_g)/\sqrt{\sigma^2_g}$;
values below 2 indicate unreliable entropy estimates.
In our model, $z_{\text{safe}} < 2$ is reached at $c \approx 10$
($\mu_g = 0.4$, $\sigma^2_g = 5.0$), limiting the numerically
verifiable range.  Beyond this, the $c^2$ growth is predicted
by the bound but cannot be confirmed computationally without
increasing $r_{\max}$.

%----------------------------------------------------------------------
\section{Numerical Evidence}
%----------------------------------------------------------------------

The following values were obtained from \texttt{diagnose\_extreme.py} with
a trained model (seed=42, cell=8, $M=50$, $n_{\text{train}}=50$, default\_gpy).
Trained parameters: $A = 0.0265$, $\lambda_0 = -1.686$.

\begin{center}
\begin{tabular}{lrrrrrl}
\hline
Image & $\mu$ & $\sigma^2$ & $\mu_g$ & $\sigma^2_g$ & angle (rad) & notes \\
\hline
$x_t$ (target) & 7.86 & 70.9 & $-1.48$ & 0.050 & 0 & natural image \\
$x^*_{\text{DA}}$ & 101 & 3{,}240 & 0.99 & 2.28 & 0.114 & DA-optimized \\
$x^*_{\text{Std}}$ & 68{,}497 & $2.85 \times 10^8$ & 1816 & $2.0 \times 10^5$ & 0.866 & Std-optimized \\
natural & $[-6, 16]$ & $[57, 267]$ & $[-1.8, -1.3]$ & $[0.04, 0.19]$ & $[0.6, 0.7]$ & pool images \\
\hline
\end{tabular}
\end{center}

\begin{center}
\begin{tabular}{lrrrrr}
\hline
Image & $\sum p(r)$ & $H_{\text{marg}}$ & $H_{\text{cond}}$ & $U_{\text{DA}}$ & var.\ ratio \\
\hline
$x_t$ & 1.000 & 0.593 & 0.583 & 0.010 & 0.000003 \\
$x^*_{\text{DA}}$ & 0.984 & 2.836 & 2.036 & 0.800 & 0.060 \\
$x^*_{\text{Std}}$ & 0.000007 & 0.00009 & 0.00009 & $\sim$0 ($U_{\text{std}} = \infty$) & 0.9999 \\
natural & $\sim$1.0 & $[0.49, 0.69]$ & $[0.49, 0.69]$ & $[0.0001, 0.006]$ & $[0.82, 1.0]$ \\
\hline
\end{tabular}
\end{center}

Key observations:
\begin{enumerate}
    \item \textbf{DA-optimized image ($\theta = 0.114$ rad):}
    After 5000 gradient ascent steps, $\mu(x^*) = 101$
    ($\approx 13 \times \mu(x_t)$; C-norm ratio $\sqrt{K(x^*,x^*)/K(x_t,x_t)} \approx 7.8$).
    The variance ratio $\sigma^2_{\text{cond}}/\sigma^2 = 0.060$ confirms
    strong conditioning ($\rho^2 \approx 0.94$, consistent with $\theta = 6.5°$).
    The log-firing rate mean $\mu_g = 0.99$ is in the useful band
    $[-5, 10]$ because $A = 0.0265$ compresses the raw GP mean.
    The utility $U_{\text{DA}} = 0.800$ is genuine.

    \item \textbf{Std-optimized image ($\theta = 0.866$ rad):}
    Standard utility diverged after 1479 steps (utility $= \infty$ from overflow).
    With $\mu_g = 1816$, the Laplace approximation catastrophically fails
    ($\sum p(r) = 7 \times 10^{-6}$). The entropy is numerically zero.
    The \texttt{nd\_utility\_new} formula involves $\exp(\mu_g + \sigma^2_g/2)$
    which overflows float32. Note the large angle ($0.87$ rad vs $0.11$ for DA):
    without conditioning to preserve alignment, the standard utility optimizer
    diverged in a different direction.

    \item \textbf{Natural images ($\theta \approx 0.65$ rad):}
    The large angle to the target gives weak conditioning (variance ratio $> 0.82$),
    resulting in tiny utility $U_{\text{DA}} < 0.006$.

    \item \textbf{Threshold for DA collapse:}
    DA utility will eventually collapse when $\mu_g > 10$ (edge of useful band).
    This requires $\mu(x^*) > (10 - \lambda_0)/A \approx 441$, i.e.\
    $c > 441/\mu(x_t) \approx 56$.
    At the current DA-optimized image with $\mu = 101$, there is
    still substantial headroom ($101/441 \approx 23\%$ of threshold).
\end{enumerate}

%----------------------------------------------------------------------
\section{Summary}
%----------------------------------------------------------------------

The divergence of unconstrained gradient ascent on DA utility has a precise
mathematical cause:

\begin{enumerate}
    \item \textbf{Proved} (Theorem~\ref{thm:proportionality}, Corollary~\ref{cor:rho}):
    For the arc-cosine kernel with $\sigma_0^2 = 0$, scaling an image
    $x^* = c \, x_t$ preserves the posterior correlation exactly at $\rho = 1$,
    regardless of $c$. The cross-kernel vector scales as $k(cx) = c \, k(x)$,
    making the posterior covariance structure proportional.
    For $\sigma_0^2 > 0$, the fractional error converges to a constant
    $O(\sigma_0^2/v_t) \approx 0.6\%$ (Remark~\ref{rem:sigma0});
    $\rho^2$ converges to $\approx 0.987$ (not 1.0), giving a residual
    variance ratio of $\approx 1.3\%$.

    \item \textbf{Proved} (Corollary~\ref{cor:moments}):
    The posterior mean scales as $\mu(cx) = c \, \mu(x)$ and the posterior
    variance as $\sigma^2(cx) = c^2 \, \sigma^2(x)$.

    \item \textbf{Proved} (Proposition~\ref{prop:gaussian}, Remark~\ref{rem:gaussian_growth}):
    Even for Gaussian observations, utility grows as $\log(c)$
    (unbounded but slow). The key issue is not unique to the Poisson model.

    \item \textbf{Not proved, numerically confirmed} ($\alpha \approx 1.9$):
    For the Poisson-GP with exponential link, utility grows as
    $U_{\text{DA}} \sim c^\alpha$ with $\alpha \approx 1.9$ in the
    Laplace-valid range ($c \leq 10$, $z_{\text{safe}} > 1.5$).
    This is consistent with the heuristic lower bound $\propto c^2$
    and much faster than the $\log(c)$ growth in the Gaussian case.
    The exponential link function amplifies GP variance into
    firing-rate uncertainty nonlinearly.

    \item \textbf{Numerically confirmed}:
    After 5000 gradient ascent steps ($\theta = 0.114$ rad, $\mu = 101$),
    the DA utility of 0.800 is genuine (Laplace approximation valid,
    $\sum p(r) = 0.984$, $\mu_g = 0.99$ in useful band).
    The optimizer found a real loophole, not a numerical artifact.
    Additionally, fixed-angle scaling experiments (test C8) confirm that
    utility growth with norm is strongly angle-dependent: effective at
    small angles ($\theta < 0.2$ rad) where conditioning is strong,
    negligible at large angles ($\theta > 0.6$ rad) where conditioning
    is weak.
\end{enumerate}

The root cause is that the arc-cosine kernel is a \emph{positively homogeneous}
function of its inputs: $K(cx, y) = c \, K(x, y)$ when $\sigma_0^2 = 0$.
This makes norm and direction completely decoupled in the GP posterior.
Gradient ascent exploits the norm degree of freedom to increase variance
(and thus entropy and utility) without sacrificing the conditioning
benefit provided by directional alignment.

Any practical use of gradient-based image optimization with this kernel
requires either a norm constraint or a normalized kernel
$\bar{K}(x,x') = K(x,x')/\sqrt{K(x,x) K(x',x')}$ that depends only on angle.

\end{document}
